\begin{table}[H]
\centering
\caption{Age-range sensitivity analysis in the 2022--2023 temporal-validation cycle. The phenotype with the highest 15--44 prevalence ratio for each endpoint is shown, then re-evaluated after including females aged 45--49 years. Phenotype assignment used the primary endpoint-excluded encoder without refitting; full-age assignment labels were aligned to the primary P0/P1/P2 labels by majority overlap in the original 15--44 subset, and reassignment agreement was checked only as a technical reproducibility control.}
\label{tab:age_sensitivity}
\scriptsize
\setlength{\tabcolsep}{2pt}
\begin{tabular*}{\textwidth}{@{\extracolsep{\fill}}>{\raggedright\arraybackslash}p{0.25\textwidth}crrrr>{\raggedright\arraybackslash}p{0.14\textwidth}@{}}
\toprule
Endpoint & Phenotype & 15--44 N & 15--44 PR & 15--49 N & 15--49 PR & Direction \\
\midrule
Contraceptive vulnerability & P1 & 2,228 & 1.33 & 2,691 & 1.26 & Same direction \\
Fertility or loss help & P2 & 239 & 1.68 & 315 & 1.35 & Same direction \\
Mistimed/unwanted pregnancy history & P2 & 239 & 2.23 & 315 & 1.98 & Same direction \\
Adverse pregnancy-history proxy & P2 & 239 & 4.11 & 315 & 3.25 & Same direction \\
Fecundity limitation/infertility & P2 & 239 & 1.47 & 315 & 1.34 & Same direction \\
\bottomrule
\end{tabular*}
\end{table}
